\section{Amplituhedron}
\label{sec:amplituhedron}

\textbf{Literature.} The amplituhedron program reformulates planar $\mathcal{N}=4$ super-Yang-Mills scattering amplitudes as the canonical form of a \emph{positive geometry}, so that locality and unitarity are derived consequences rather than starting axioms. The sources below goes from the BCFW recursion of Section~\ref{sec:bcfw-recursion} to the present.

\begin{itemize}
    \item \textbf{Foundations.} \autocite{Britto:2005fq} is the BCFW recursion itself (Section~\ref{sec:bcfw-recursion}), whose individual terms the amplituhedron geometrizes. \autocite{hooftPlanarDiagram1974} introduced the large-$N$ planar limit and the $1/N$ expansion (Section~\ref{sec:color-ordering-planarity}); positivity and the amplituhedron are developed for planar theories.

    \item \textbf{On-shell diagrams and the positive Grassmannian.} \autocite{Arkani-Hamed:2012zlh} gives on-shell diagrams (made by gluing three-point amplitudes) a geometric meaning: each diagram labels a cell of the positive Grassmannian $G_+(k,n)$. In this language, BCFW steps become simple adjacent swaps in the diagram's permutation data, and Yangian symmetry appears as the freedom to move within the same positive geometry without changing the amplitude.
    This is the technical content behind the on-shell-diagram and positive-Grassmannian formulation of the amplitude.

    \item \textbf{The all-loop integrand.} \autocite{Arkani-Hamed:2010zjl} promotes BCFW to a recursion for the all-loop \emph{integrand} of planar $\mathcal{N}=4$ SYM in momentum-twistor space (also discussed in Section~\ref{sec:loop-level-recursion}).

    \item \textbf{The amplituhedron.} \autocite{Arkani-Hamed:2013jha} defines the amplituhedron $\mathcal{A}_{n,k,L}$, a generalization of the positive Grassmannian whose canonical form (its ``volume'') yields the amplitude, with locality and unitarity emerging from the positive geometry rather than imposed. \autocite{Damgaard:2019ztj} constructs the \emph{momentum amplituhedron}, the analogous positive geometry directly in spinor-helicity space (Section~\ref{sec:spin-helicity-formalism}) rather than momentum-twistor space, for tree-level amplitudes.

    \item \textbf{Recent mathematical proofs.} \autocite{Even-Zohar:2021sec} proves the Arkani-Hamed--Trnka conjecture that the $m=4$ amplituhedron decomposes into images of BCFW positroid cells (the ``BCFW triangulation''), and \autocite{Even-Zohar:2025ydi} proves the associated tiling and cluster-adjacency conjectures. Together they make rigorous the BCFW--amplituhedron correspondence: every way of iterating the BCFW recursion yields a tiling of one and the same geometry.

\end{itemize}